\section{Conclusion}
\label{sec:conclusion}

This project built and evaluated eight NLP models from scratch, covering the full
spectrum from simple probabilistic tables to attention-based neural architectures.
The key findings are:

\begin{enumerate}
    \item \textbf{LSTM gating resolves vanishing gradients in practice.}
          The character-level LSTM achieved a final loss of 0.2943 versus 0.6689
          for the character-level RNN — a 56\% reduction — despite identical
          architecture depth and training configuration. The difference is
          attributable entirely to the forget, input, and output gates.

    \item \textbf{Tokenization granularity interacts with dataset size.}
          Character-level models outperformed word-level models on this corpus
          because characters repeat more predictably across 12,000 words. Word-level
          models require significantly more data to see enough unique word-context
          combinations for reliable next-word prediction.

    \item \textbf{Architecture alone does not determine convergence speed.}
          The Transformer chatbot trained faster per epoch but did not converge
          within 30 epochs on 2,000 QA pairs, while the LSTM chatbot converged to
          a loss of 0.10 on the same data and time budget. Training stability
          (teacher forcing), parameter count, and learning-rate schedule all
          interact with architecture choice.

    \item \textbf{Repetition collapse is a concrete failure mode of under-trained
          sequence models.} The Transformer's collapse into repetitive phrase loops
          illustrates why beam search, top-$k$ sampling, or nucleus sampling are
          important additions to any deployed text-generation system.
\end{enumerate}

\subsection*{Limitations and Future Work}

\begin{itemize}
    \item All models are trained on procedurally generated text. Real corpora
          (e.g., Wikipedia, Common Crawl) would expose models to genuine linguistic
          diversity and allow more meaningful qualitative evaluation.
    \item The Transformer chatbot would benefit from training for 100–200 epochs,
          or from replacing the warmup schedule with a cyclic LR to converge
          faster on small data.
    \item Beam search decoding (rather than greedy or temperature sampling) would
          likely improve the coherence of all generative models.
    \item Quantitative evaluation metrics such as BLEU, perplexity, or BERTScore
          would complement the current loss-based comparison.
\end{itemize}
